% !TeX root = ../papers/nc-p3.tex
\section{安定な非正則軌跡の既知十五次元成分への大域的同定}
\label{sec:binary-global-component}
\begin{theorem}\label{thm:binary-global-component}
定理\ref{thm:binary-component-crossing}の15次元成分は
定理\ref{thm:full-fifteen-component}の既約成分と同一である。
非自明な二捩れ楕円鉛筆の13次元軌跡と、定理\ref{thm:binary-linear-complex}の
有限非正則窓の全てがこの成分に属する。両軌跡の一般点は異なり、前者は両対称で正則、
後者は片側対称で非正則である。

この成分の一般の単一二元辺を持つ非正則窓では全窓の軌跡は滑らかで、
片側対称因子への横断変数を $a$ とすると反対行列式は単元倍の $a^6$ である。
従って前節までの残余階数15因子はこの片側対称な非正則13次元軌跡を含まない。
一般二元点の近くで二つの軌跡が片側対称因子上に切る13次元部分も互いに異なる。
\end{theorem}
\begin{proof}[一つの印からの成分の照合]
線型複体の非退化部分において、印を選ぶ有限被約開集合は空でない。
以下、滑らかな判別三次を持つ印が一般の窓に存在することを示す。三次形式を
$q=1-x+x^3$ とし、$L(q)=(2+3x+x^2)q$ となる有理標本を
\path{computations/verify_binary_global_component.py}で取る。
指定点で六つの双一次式のアフィンJacobi行列は可逆なので、
この印は辺と共に局所エタールに動き、入射の像はパラメータ空間の開集合を含む。
$D\in S$ 側の座標は $H_{5,3}(D)q=0$ の唯一の核線で、標本での非零小行列により近傍にも戻せる。
一般に有限Frobeniusであることは前定理による。従って印の開条件の検査を、この有理標本で行えばよい。

$U,Q,S,Q$ 順に回転して印を先頭の座標とし、反対評価の二行列を単位に戻す。
この片側対称表示では共通の非退化対称計量で $G,F$ は自己随伴である。
直交基底で $P=Q$ かつ $G,F$ が対称の表示にできるが、判別三次を求めるだけなら
計量の平方根を取る必要はない。$[G_i,F_j]$ の三次元線型包の三座標を取れば、
対応する三行三列一次式行列の行列式が同じ平面曲線を与える。
素数1009で三つの偏微分は共通射影零点を持たない。
従って標数零の印付き開部分にも滑らかな判別三次があり、
定理\ref{thm:marked-smooth-discriminant-components}からその窓は既知の15次元成分に入る。

同じ標本で $Q$ の反対交換の二指標の接次元の上界は $60,16$ である。
一方、一般の単一二元窓は有限かつ非両対称なので全成分のアフィン次元が76以上である。
よって一般点は滑らかな76次元で、成分は一つしかない。
この点は二元境界から得た15次元成分にも入る。実際既約な単一二元辺の軌跡は
一般両対称二元点を含み、そこでの二つの局所成分のうち、両対称でない一般点を受けられるのは
15次元のものだけである。以上で二成分の大域的な一致を得る。
全ての許容有限パラメータも、既約な入射からの連続なテンソルの構成により同成分に入る。
\end{proof}
\begin{proof}[片側対称因子の横断]
一つの交換の正指標の160変数には一切断の100本のSchur式だけが課されるので、
正指標の接次元は少なくとも60であり、一般点ではちょうど60である。
軌道の正負指標次元はそれぞれ46と15である。従って軌道横断の15変数で
固定部分は滑らかな14次元、負指標は一変数 $a$ のみである。
交換による同定は有限グラフ上で唯一なので、前節の局所議論と同じく基底変更による飽和もこの固定条件となる。
反対切断で交代六行の $a$ をくくった行列の、$a=0$ での可逆性を調べる。
正十行は階数10であるから、残る六次元核から交代六行への負接による写像を取ればよい。
\path{computations/verify_binary_linear_complex.py}の標本でこの階数は6である。
負接の一般次元16は上の下界と標本のJacobi階数の双方で確認できるため、
核とその変分を良い素数から標数零へ持ち上げられる。可逆性は既約な族の稠密開集合で成り立つ。
従ってこの点には片側対称因子以外の反対行列式の局所成分は来ない。
\end{proof}
